% Problem setting and requirements. Sources: AGENT/research_content.md §1.5, §3 (intended logic), docs/heldout_design.md
% §1-§2 and §14 (as built), docs/baselines/authoring_spec.md, the SRB Mission A task (brief text). Where the build and
% research_content.md disagree, the build is followed and the disagreement is noted in a comment.
% Preamble needs \usepackage{enumitem} for the R1-R6 labels.
\providecommand{\TBD}[1]{\textcolor{red}{[TBD: #1]}}

\section{Problem setting and requirements}
\label{sec:requirements}

Lunar surface infrastructure will run for long periods without crew, and robots are expected to keep it running. They
will reset tripped power units, replace failed modules and move spares between stands while the ground team is far
away \citep{lewis2022uncrewed}. Mission control states such work as an objective. A representative objective in this
study reads: ``PDU-2 reported a fault. Inspect the base's serviced units, identify the failed unit, restore it if within
your authority, verify the result, and return the rover to the dock with the arm stowed.'' The objective names neither
the cause nor the steps. The robot has to find both.
% Brief text: space_robotics_bench/srb/tasks/mobile_manipulation/lunar_base_mission_a/task.py mission_brief(),
% L0 instance faulted_unit=pdu_2 (docs/heldout_design.md §14).

Mission control cannot sit inside every decision. Operational one-way delays for lunar missions range from 3 to 14\,s
\citep{landon2025delay}. From a south-polar site, Earth is above the horizon for only about half of each sidereal month
\citep{nasa2023lcn}. For days at a time, a request to the ground is stored and forwarded, not answered. Onboard autonomy
already absorbs this for the situations its designers modelled. Plan executives run authored procedures with conditions,
loops and recovery branches \citep{estlin2006plexil}. Model-based planners replan within an authored domain
\citep{bedraxweiss2004europa,rabideau2025cadre}. This paper concerns the situations whose response nobody authored.
% "3 to 14 s": Landon et al. 2025 (research_content.md §1.5.1). "about half of each sidereal month": NASA LCN white
% paper 2023, as cited in space_context/earth_visibility.py and docs/heldout_design.md §2.

The comparison fixes everything except the mission-level system. Every method drives the same robot through the same
closed set of capabilities behind one typed contract. All methods share the sensors, the vision-language model and the
low-level safety kernel, and each call returns the same raw evidence to whichever method made it. One external evaluator
scores every run. It reads simulator truth, and no method does. What differs is how a method turns the objective and the
evidence into calls. The baselines are not given the proposed runtime's mechanisms. Each encodes mission and authority
constraints in its own artifacts, as its designers intended. The comparison therefore holds the problem and the robot
fixed, and lets the mechanisms differ.
% research_content.md §3.2 "Method boundary (revised 2026-09-17)"; the evaluator is the only reader of simulator truth
% (evaluation/violations.py docstring; openral_ext/capability_server/node.py docstring: "No simulator truth").

A situation is \emph{covered} for a baseline when its frozen artifacts contain the procedure, recovery branch or domain
relation that a valid solution needs. A situation is \emph{held-out} when five conditions hold together. Every
capability a solution needs exists and is available to every method. The relevant physical state is observable with the
shared sensors. The frozen baseline artifacts encode no decomposition, relation or composition for the situation. The
proposed runtime receives no hidden solution either. It gets the generic objective, the capability contracts, the
current evidence and the frozen authority configuration. No artifact of any method changes once the situation is
revealed, and a run that could continue only after such a change is recorded as requiring ground re-authoring.
Held-out situations are graded by what the frozen baseline model lacks. The ladder runs from a missing cause or relation
to a missing composition across mission families, and then to a goal conflict that needs ground authority
(Section~\ref{sec:setup-missions}). The grade depends on the baseline's artifacts alone, so the ladder does not depend on
how hard a case is for the proposed runtime.
% research_content.md §3.3 (the five conditions); docs/heldout_design.md §1 rule 5 (levels by what the frozen baseline
% model lacks).

The setting yields three hypotheses. H1 predicts that, on the same robot, the proposed runtime completes held-out
situations without ground re-authoring where the frozen baselines stop. H2 predicts that, when a capability is added
after the freeze, the proposed runtime needs less mission-specific engineering to reach validated support, with
comparable held-out success. H3 predicts that this autonomy does not come from crossing hard boundaries. Under H3,
authority and invariant boundaries hold and no unverified effect enters the mission record. No completion is declared
that the evaluator does not confirm, and the runtime contains no task order or recovery for any situation.
% research_content.md §1.7 (H1-H3). H2 is stated as registered; Section results-exp2 reports that active time was
% similar and the kind of change differed, which the results section argues, not this one.

Six requirements follow for any system that gives a foundation model the mission strategy. The first comes from the
question under test, and the other five from lunar operations.
\begin{enumerate}[label=R\arabic*]
  \item \label{req:no-solution} The runtime pre-engineers boundaries and evidence, not solutions. It may encode
  capabilities, facts, verification and authority. It may not encode a task order, an observation policy or a recovery
  for a situation. A held-out situation whose response is hidden in the runtime is in fact covered, and H1 would then
  test nothing.
  \item \label{req:facts} Mission truth lives in an authoritative record that only evidence advances. Flight computers
  are limited in power and radiation tolerance \citep{nasa2024hpsc}, so the strongest model may run behind a link that
  closes. The record must survive a restart or a change of reasoning service. A model's belief must never become a
  fact.
  \item \label{req:verification} Every physical effect is verified independently of the model and of the skill that
  produced it. With the ground out of the loop, nobody watches each action. A skill's success flag says the skill
  finished, and a model's report says what the model believes. Neither shows that the world changed.
  \item \label{req:authority} Authority and invariant boundaries are enforced outside the model. Flight rules reserve
  some operations for mission control, such as removing equipment from operating infrastructure \citep{cfr1214403}. A
  boundary that holds only when the model complies is advice.
  \item \label{req:fdir} Failures are detected and isolated onboard. They are escalated to the next level of authority
  when onboard handling cannot resolve them \citep{ecss2008operability}. The escalation crosses a delayed link and may
  wait for the next pass \citep{cerf2007dtn}. The robot must not block on it, and it must know the timing of the link.
  \item \label{req:trace} Mission control can reconstruct what the robot knew, what it did and why. More onboard
  autonomy moves the ground's role from commanding each step to reconstructing the vehicle's decisions afterwards
  \citep{castano2022ops}.
\end{enumerate}
% Disagreement with research_content.md §3.4: it lists five requirements, including "semantic observation" (§3.4.2) and
% "space-native operational context and distributed intelligence" with onboard / edge / Earth tiers (§3.4.3). As built,
% semantic observation is one query of the shared capability set (observe, answered by the shared VLM), and the
% space-native context is the Direct-with-Earth link state, the units' telemetry and mission control's decision time;
% no edge or onboard reasoning tier exists (the reasoner is one remote model, configs/space_runtime.yaml). The
% distributed-intelligence argument survives here only as R2's restart and endpoint-change requirement, which the
% record supports (space_runtime/facts.py reasoner_history; agent.py resume option) and ablation J3 tests.
% The FDIR and trace requirements are split out of research_content.md §3.4.4/§3.4.5 to match the user's framing
% (five engineering disciplines plus model-owned strategy).

Requirement R1 separates the proposed runtime from a mission executive. R2 to R6 are disciplines that space operations
already require of onboard autonomy. Section~\ref{sec:method} builds each of R2 to R6 as a mechanism outside the model,
and it keeps R1 by leaving every choice of action to the model.
